\PassOptionsToPackage{numbers}{natbib}
\documentclass{article}
\usepackage{iclr2024_conference}
\usepackage{luatexja-fontspec}
\usepackage{hyperref}
\usepackage{url}
\usepackage{booktabs}
\usepackage{amsmath}
\usepackage{amsfonts}
\usepackage{graphicx}
\graphicspath{{../}}

\InputIfFileExists{values.tex}{}{}
\providecommand{\airasval}[1]{\textbf{??airasval:\detokenize{#1}??}}
\providecommand{\unverified}[1]{#1}
\providecommand{\airasrecordlink}[1]{#1}

\title{SciGym における自前 serving の NVIDIA Nemotron 2 種の事前登録評価: 論文 Table 1 と同じ条件で small を測り、large にも広げる}
\author{AIRAS}

\begin{document}
\maketitle

\begin{abstract}
SciGym \citep{duan-2025-measuring} は、反応をすべて取り除いた SBML モデルを LLM エージェントに渡し、摂動実験とコード実行を繰り返して反応機構を復元させるベンチマークである。論文の Table 1 は 2025 年春の 3 社のフロンティアモデル 6 つを SciGym-small 137 件で各 1 回走らせた平均値を報告している。本研究は、公開コードの Controller と評価器をそのまま用い、論文と同じ条件（137 件、最大 20 反復、temperature 1.0）で、NVIDIA の open-weight モデル 2 種（nemotron-3.5-lightning、30B で活性 3B の MoE と、nemotron-3-ultra、550B で活性 55B の MoE）を評価する。両モデルは外部の推論 API ではなく、RIKYU の GB200 ノード上に NIM コンテナで自前 serving する。あわせて、論文が評価していない large split 213 件も同じ条件で走らせる。実験前に、small と large のそれぞれで 2 モデルの平均 Simulation Trajectory Error（STE）が論文 Table 1 の最良行（Gemini-2.5-Pro の 0.3212）以下であることを反証線として登録する。Reaction Matching Score（RMS）と Network Topology Score（NTS）は同じ実行から表として報告する。
\end{abstract}

\section{はじめに}
\label{sec:intro}
LLM に科学的発見のサイクル全体（実験の設計、結果の解析、機構の提案）を行わせて測るベンチマークとして、SciGym \citep{duan-2025-measuring} が提案された。BioModels 由来の SBML モデルから反応をすべて取り除いたものをエージェントに渡し、エージェントは初期濃度を変えた摂動実験を要求し、返ってきた時系列を Python で解析し、最終的な SBML を提出する。提出モデルは真のモデルと、軌道誤差（STE）と反応の一致度（RMS）で比較される。

論文の主結果である Table 1 は、3 社の pro / mini 計 6 モデルを SciGym-small の 137 件で評価した平均値である \cite[s1.p2, s1.p6]{duan-2025-measuring}。評価されたのは 2025 年春の版（Claude-3.7-Sonnet と 3.5-Haiku、Gemini-2.5-Pro と Flash のプレビュー版、GPT-4.1 と 4.1-mini）である \cite[s1.p1]{duan-2025-measuring}。AIRAS の先行研究は、RIKYU の推論 API で提供される open-weight モデル 6 種を同じプロトコルで評価し \cite[s3.p1]{scigymrikyuopenmodels-2026-scigym}、kimi-k2.6 の平均 STE が \unverified{0.2100} \cite[s3.p2]{scigymrikyuopenmodels-2026-scigym}、kimi-k3 が \unverified{0.2251} \cite[s3.p3]{scigymrikyuopenmodels-2026-scigym} と、論文 Table 1 の最良行 Gemini-2.5-Pro の \unverified{0.3212} \cite[s1.p5]{duan-2025-measuring} を下回ることを示した。

これらはいずれも、第三者が運用する推論 API を呼ぶ形の評価であった。本研究は、重みが公開されているモデルを自分の計算機で serving する場合に同じ課題がどこまで解けるかを、NVIDIA の Nemotron 系列 2 種で測る。lightning は 1 GPU で動く活性 3B の軽量モデル、ultra は 4 GPU を要する活性 55B の大型モデルであり、同じ系列の中で規模の効果も見える。論文が計算資源の制約から評価しなかった large split 213 件 \cite[s1.p2]{duan-2025-measuring} にも同じ設定で広げる。判定は AIRAS の事前登録の枠組みに従い、仮説・判定基準・予測区間・実行条件を実験前に記録へ凍結し、実験後は計測値だけを流し込む。

\section{関連研究}
\label{sec:related}
SciGym \citep{duan-2025-measuring} は、単一スキルを測る既存のベンチマークと異なり、エージェントが自ら摂動実験を設計してその結果を受け取る閉ループを評価する。評価指標は STE、RMS（modifier まで要求する厳格版もある）、NTS の 3 つで、論文は pro モデルが mini モデルを一貫して上回り、Gemini-Pro が最良であったと報告している \cite[s1.p6]{duan-2025-measuring}。

open-weight モデルの評価では kimi-k2.6 が \unverified{0.2100} \cite[s3.p2]{scigymrikyuopenmodels-2026-scigym}、kimi-k3 が \unverified{0.2251} \cite[s3.p3]{scigymrikyuopenmodels-2026-scigym}、deepseek-v4.1-flash が \unverified{0.2405} \cite[s3.p4]{scigymrikyuopenmodels-2026-scigym} に達し、27B〜35B 級は qwen3.8-27b の \unverified{0.2479} \cite[s3.p5]{scigymrikyuopenmodels-2026-scigym} と qwen3.6-35b の \unverified{0.4571} \cite[s3.p6]{scigymrikyuopenmodels-2026-scigym} に分かれた。

\section{仮説と予測}
\label{sec:hypothesis}
以下は記録（record.json）から生成された仮説・主張・判定基準・予測区間の一覧である。実験後は Observed と Verdict が計測値から埋められる。

% AUTO-GENERATED by the update_record tool. DO NOT EDIT.
% verify_paper regenerates this file from record.json and the run outputs
% and fails on any difference, so manual edits always surface.
\noindent\textbf{Hypothesis 1.} NVIDIA の open-weight モデル nemotron-3.5-lightning（30B、活性 3B の MoE）と nemotron-3-ultra（550B、活性 55B の MoE）を RIKYU の GPU ノード上の NIM サーバー（OpenAI 互換）で自前 serving し、SciGym 公開コードの Controller をそのまま用いて論文と同じ条件（SciGym-small 137 件、最大 20 反復、temperature 1.0）で解かせたとき、いずれも平均 Simulation Trajectory Error（STE）が論文 Table 1 の最良行（Gemini-2.5-Pro の 0.3212）以下である。~\cite[s1.p1, s1.p5, s1.p6]{duan-2025-measuring}, \cite[s3.p2]{scigymrikyuopenmodels-2026-scigym}
\begin{enumerate}
\item[\textbf{Claim 1}] nemotron-3.5-lightning（nvidia/nemotron-3.5-lightning）で SciGym-small 137 件を最大 20 反復で解かせたときの平均 STE は、論文 Table 1 の最良行（Gemini-2.5-Pro）の 0.3212 以下である。~\cite[s1.p5]{duan-2025-measuring}
  \emph{Rationale:} 活性 3B の軽量 MoE が論文の最良行に届けば、h1 の lightning の部分を担う。
  \emph{Criterion:} \texttt{\detokenize{nemotron-3.5-lightning-small}}.\texttt{\detokenize{trajectory_smape}} $-$ 0.3212 $\leq 0$.
  \emph{Prediction:} $[-0.07, 0.18]$ (同条件の 27B〜35B の open-weight モデルは 0.248〜0.457（s3.p5、s3.p6）。活性 3B の軽量モデルはこの幅に収まると見る).
  \emph{Observed:} 0.257174.
  \emph{Verdict:} refuted.
\item[\textbf{Claim 2}] nemotron-3-ultra（nvidia/nemotron-3-ultra-550b-a55b）で SciGym-small 137 件を最大 20 反復で解かせたときの平均 STE は、論文 Table 1 の最良行（Gemini-2.5-Pro）の 0.3212 以下である。~\cite[s1.p5]{duan-2025-measuring}
  \emph{Rationale:} 活性 55B の大型 MoE が論文の最良行に届けば、h1 の ultra の部分を担う。
  \emph{Criterion:} \texttt{\detokenize{nemotron-3-ultra-small}}.\texttt{\detokenize{trajectory_smape}} $-$ 0.3212 $\leq 0$.
  \emph{Prediction:} $[-0.22, -0.02]$ (同条件の大型 open-weight MoE（kimi-k2.6、kimi-k3、deepseek-v4.1-flash）は 0.21〜0.24（s3.p2〜s3.p4）。ultra は同等かやや良いと見る).
  \emph{Observed:} -0.0376069.
  \emph{Verdict:} supported.
\end{enumerate}
\noindent\emph{Assumed, so that the claims together imply Hypothesis 1:}
\begin{itemize}
\item temperature 1.0 での 1 回実行の平均 STE（137 件平均）の実行間変動が判定に対して十分小さいこと。先行する Table 1 再現では件ごとの STE の標準偏差はおよそ 0.27、137 件平均の標準誤差はおよそ 0.02 であった（c1、c2）
\item NIM コンテナ（vLLM）で serving したモデルが、提供元の推奨設定で呼んだ場合と同じ能力で応答すること。思考（reasoning）はモデルの既定のままとし、思考の文が本文に含まれて返る場合もそのまま Controller に渡す（c1、c2）
\item aarch64 向けの libroadrunner 2.7.0 と antimony 2.14.0 が、論文環境の libroadrunner 2.8.0 と同じ軌道と評価値を与えること（c1、c2）
\item 平均 STE が「論文の最良行に届くか」を代表すること。RMS（modifier あり / なし）と NTS は同じ実行から表として報告するが、判定基準には用いない（c1、c2）
\item 文脈長（262144 トークン）や時間切れで提出に至らなかった件は不完全モデルの採点値で平均に含める（c1、c2）
\item NIM サーバーの job が計算機の 4 日上限で打ち切られても、別ノードで立ち直したサーバーに実験側が接続し直し、反復の途中で条件が変わらないこと（c1、c2）
\end{itemize}
\noindent\textbf{Hypothesis 2.} 同じ 2 モデルを、論文が評価していない SciGym-large 213 件（反応数の多い系）で同じ条件（最大 20 反復、temperature 1.0）で解かせたとき、いずれも平均 STE が論文 Table 1 の最良行（small での Gemini-2.5-Pro の 0.3212）以下である。~\cite[s1.p2, s1.p5]{duan-2025-measuring}, \cite[s3.p2]{scigymrikyuopenmodels-2026-scigym}
\begin{enumerate}
\item[\textbf{Claim 3}] nemotron-3.5-lightning で SciGym-large 213 件を最大 20 反復で解かせたときの平均 STE は、論文 Table 1 の最良行（Gemini-2.5-Pro）の 0.3212 以下である。~\cite[s1.p5]{duan-2025-measuring}
  \emph{Rationale:} 軽量モデルが大規模な系でも論文の最良行に届けば、h2 の lightning の部分を担う。
  \emph{Criterion:} \texttt{\detokenize{nemotron-3.5-lightning-large}}.\texttt{\detokenize{trajectory_smape}} $-$ 0.3212 $\leq 0$.
  \emph{Prediction:} $[0.03, 0.28]$ (large には先行値が無い。small での予測（0.25〜0.50）に、反応数の多い系での悪化を 0.1 程度見込む).
  \emph{Observed:} 0.275656.
  \emph{Verdict:} refuted.
\item[\textbf{Claim 4}] nemotron-3-ultra で SciGym-large 213 件を最大 20 反復で解かせたときの平均 STE は、論文 Table 1 の最良行（Gemini-2.5-Pro）の 0.3212 以下である。~\cite[s1.p5]{duan-2025-measuring}
  \emph{Rationale:} 大型モデルが大規模な系でも論文の最良行に届けば、h2 の ultra の部分を担う。
  \emph{Criterion:} \texttt{\detokenize{nemotron-3-ultra-large}}.\texttt{\detokenize{trajectory_smape}} $-$ 0.3212 $\leq 0$.
  \emph{Prediction:} $[-0.12, 0.13]$ (large には先行値が無い。small での予測（0.10〜0.30）に、反応数の多い系での悪化を 0.1 程度見込む).
  \emph{Observed:} 0.111574.
  \emph{Verdict:} refuted.
\end{enumerate}
\noindent\emph{Assumed, so that the claims together imply Hypothesis 2:}
\begin{itemize}
\item large split には論文の公表値が無いので、論文の small での最良値を絶対水準の基準とする。small との難易度の差は表で並べて示すが、判定には用いない（c3、c4）
\item 213 件平均の実行間変動が判定に対して十分小さいこと。件ごとの標準偏差を 0.27 とすれば平均の標準誤差はおよそ 0.02（c3、c4）
\item 評価層 scigym\_large（airas-eval v0.13.1）が公式リリースの large split 213 件を sha256 で固定し、small と同じ公式 Evaluator で採点すること（c3、c4）
\item 文脈長（262144 トークン）や時間切れで提出に至らなかった件は不完全モデルの採点値で平均に含める（c3、c4）
\end{itemize}


\paragraph{判定基準と予測区間の根拠.}
各主張の反証線は「（当該モデルの平均 STE）$-$ \unverified{0.3212} $\le 0$」である。\unverified{0.3212} は論文 Table 1 の最良行 Gemini-2.5-Pro の STE \cite[s1.p5]{duan-2025-measuring} であり、small（仮説 1）と large（仮説 2）の 4 主張すべてに同じ基準を置く。large には論文の公表値が無いので、small での最良行を絶対水準の基準とする。予測区間は、small の lightning で $[-0.07, 0.18]$（同条件の 27B〜35B 級が 0.25 と 0.46 に分かれたこと \cite[s3.p5, s3.p6]{scigymrikyuopenmodels-2026-scigym} による）、small の ultra で $[-0.22, -0.02]$（同条件の大型 open-weight MoE が 0.21〜0.24 \cite[s3.p2, s3.p3, s3.p4]{scigymrikyuopenmodels-2026-scigym} であったことによる）とし、large は先行値が無いので small の予測に 0.1 程度の悪化を見込んで lightning $[0.03, 0.28]$、ultra $[-0.12, 0.13]$ とする。1 件の STE の標準偏差を 0.27 程度と見積もると 137 件平均の標準誤差は約 0.02 で、差が 0.02 未満のときは結論が揺らぎうることを前提として記録する。

RMS（modifier あり / なし）と NTS も同じ実行から表として報告するが、1 つの主張に置ける判定基準は 1 つであるため、判定には用いない。lightning と ultra の差、small と large の難易度差も表で示すにとどめ、主張にはしない。

\section{方法}
\label{sec:method}
\paragraph{タスク.}
各インスタンスは真の SBML（\texttt{truth.xml}）、反応を取り除いた不完全な SBML（\texttt{partial.xml}）、シミュレーション条件（\texttt{truth.sedml}）からなる。エージェントは不完全な SBML を受け取り、各ターンで実験（初期濃度の変更を指定し、真のモデルの時系列を得る）、コード実行（numpy / pandas / scipy / libsbml と、仮説モデルを走らせる \texttt{simulate} 関数が使える）、提出のいずれかを markdown 形式で返す。提出は任意の時点で行え、反復数の上限は 20 である \cite[s1.p3]{duan-2025-measuring}。提出した SBML が無効な場合は追加のデバッグ反復が許され \cite[s1.p4]{duan-2025-measuring}、それでも無効なら不完全 SBML のまま採点される。

\paragraph{実行系.}
ループ、プロンプト、実験の実行はすべて公開コード \citep{h4duan-2025-scigym} の \texttt{Controller} をそのまま用いる。本研究が書くのは、(1) OpenAI 互換 API でモデルを呼ぶ \texttt{LLM} 実装（公式の GPT 実装と同じく、システムプロンプトと全履歴を毎回送る）、(2) インスタンスごとに 1 プロセスで Controller を起動し、複数件を並列に走らせる駆動部、(3) 各インスタンスの提出モデルと公式データを評価層の入力ファイルに書き出す部分だけである。公式コードでは履歴がクラス属性として定義されており 1 プロセスで複数件を走らせると履歴が持ち越されるが、1 件 1 プロセスとすることでこの問題を避ける。LLM が書いた解析コードや提出モデルの評価が ODE ソルバーの C ライブラリ内で止まると公式の 3 分制限（\texttt{signal.alarm}）が効かないので、駆動部は出力が 60 分止まった試行を打ち切って反復 0 からやり直す（最大 3 試行）。

\paragraph{serving.}
2 モデルは NVIDIA NIM のコンテナ（lightning は nemotron-3.5-lightning-30b-a3b 2.0.9、ultra は nemotron-3-ultra-550b-a55b 2.0.11、NVFP4 量子化）を Apptainer で RIKYU の GB200 ノードに立て、OpenAI 互換のエンドポイントとして呼ぶ。lightning は 1 GPU、ultra は 4 GPU の Slurm job で、いずれも文脈長は 262144 トークンである。サーバーの job は計算機の 4 日上限で打ち切られうるので、実験側はサーバーの接続先を共有ディスク上のファイルから読み、接続断のたびに読み直して別ノードで立ち直ったサーバーに繋ぎ直す。思考（reasoning）はモデルの既定のままとし、lightning のように思考の文が本文に含まれて返る場合もそのまま Controller に渡す。

\paragraph{指標.}
採点は評価層 airas-eval のベンチマークパック \texttt{scigym\_small}（137 件）と \texttt{scigym\_large}（213 件）が行う。パックは公式リリースの各件（truth / partial / sedml）を sha256 で固定し、公開コードの \texttt{Evaluator}（コミット 88a7b93）で提出モデルを採点して、STE（\texttt{trajectory\_smape}）、RMS の適合率・再現率・F1（modifier なし \texttt{reaction\_*}、あり \texttt{reaction\_*\_with\_modifiers}）、NTS の F1（\texttt{topology\_f1}）を件の単純平均で返す。提出 SBML が無い、または無効な件は、公式評価どおり不完全 SBML の採点値で平均に含める。

\section{実験設計}
\label{sec:design}
\paragraph{モデル.}
nvidia/nemotron-3.5-lightning（30B、活性 3B の MoE）と nvidia/nemotron-3-ultra-550b-a55b（550B、活性 55B の MoE）の 2 種。いずれも重みが公開されており、NIM の推奨設定（generation config の temperature 1.0、top-p 0.95）で serving する。

\paragraph{データ.}
SciGym の公式リリース（Hugging Face \texttt{h4duan/scigym-sbml}）のうち、論文が評価対象とした small split 137 件 \cite[s1.p2]{duan-2025-measuring} に加え、論文では評価されていない large split 213 件を用いる（いずれも RIKYU 上のコピー）。評価層は各件の内容を公式リリースの sha256 と照合する。1 run は 1 モデル × 1 split で、計 4 run である。

\paragraph{設定.}
\texttt{max\_iterations} は論文どおり 20 \cite[s1.p3]{duan-2025-measuring}（公開 config の既定値は 5 \cite[s2.p2]{h4duan-2025-scigym}）、\texttt{eval\_debug\_rounds} は論文本文の 3 \cite[s1.p4]{duan-2025-measuring}（公開 config の既定値は 5 \cite[s2.p3]{h4duan-2025-scigym}）、\texttt{temperature} は公開 config の 1.0 \cite[s2.p1]{h4duan-2025-scigym}、摂動は初期濃度の変更のみとする。応答のトークン上限は 32768 とし、思考で使い切った応答は文脈長に収まる範囲で上限を倍にして呼び直す。文脈長を超えて提出に至らなかった件は不完全 SBML の採点値で平均に含め、件数を表に出す。各モデル、各インスタンスを 1 回ずつ実行する。並列は lightning で 48 件、生成の遅い ultra で 16 件とし、1 件 1 試行の上限は 4 時間とする。

\paragraph{計算資源と環境.}
実行は Seyval を通じて RIKYU（Slurm、GB200 ノード、aarch64）で行い、LLM は同じ計算機上の NIM サーバーを呼ぶ。実験側は GPU を使わないが、資源の単位が GPU であるため 1 GPU を要求する。環境は Dockerfile と \texttt{uv.lock} で固定する。aarch64 向けの wheel が無いため、libroadrunner は公式の 2.8.0 ではなく 2.7.0、antimony は 2.14.0 を用い、rrplugins と phrasedml は除外する。この差でシミュレーション結果がわずかに変わる可能性があり、仮説の前提として記録する。

\paragraph{段階.}
sanity は small の先頭 1 件を 2 反復で走らせて配管を確かめる。主張の判定は full（small 137 件、large 213 件）の結果だけで行い、sanity の結果は記録に取り込まない。

\paragraph{run と主張の対応.}
\texttt{nemotron-3.5-lightning-small} $\to$ Claim 1、\texttt{nemotron-3-ultra-small} $\to$ Claim 2、\texttt{nemotron-3.5-lightning-large} $\to$ Claim 3、\texttt{nemotron-3-ultra-large} $\to$ Claim 4。

% ===== airas: post-experiment sections, filled once runs exist =====

\bibliographystyle{iclr2024_conference}
\bibliography{references}

\end{document}
